\chapter{Project 4: Candle}

\section{Overview}
Now that you have a handle on how to work with LEDs and buttons, let's look at another kind of component
that you might use, a microphone. Over the course of this project, you will:
\begin{itemize}
    \item Connect a microphone board to your microcontroller
    \item Use MicroPython to poll the microphone over and over and use its value to control the brightness of an LED
    \item Connect a screen to your microcontroller and show some data on it
\end{itemize}
At the end of this project, your microcontroller should run a MicroPython program that prints out the current
temperature and humidity of the room you're sitting in. Let's get started!
\begin{figure}[H]
\centering
    \includegraphics[width=.6\linewidth]{project_4/success!.jpg}
    \caption{The end result should look something like this}
\end{figure}

\pagebreak

\section{Directions}

\subsection{Creating the circuit}
Using jumper cables, you will be assembling a circuit between your microcontroller, your breadboard, and the small
microphone board and OLED screen included in your kit.

\subsubsection{Remove previous components}
Before beginning, remove any components from prior chapters including LEDs, buttons, and wires. You may leave the
microcontroller attached to the breadboard.

\subsubsection{Attach the microcontroller to the breadboard}
If it's not already, carefully insert the pins at the bottom of your microcontroller into the breadboard. Refer back
to \ref{pinout} for pin labels. When placing the board into the breadboard, make sure that the microcontroller is oriented such that:
\begin{itemize}
    \item The pin labeled \textbf{5V} is inserted in hole at \textbf{Column H, Row 1} of the breadboard (or \textbf{H1}, for short)
    \item The pin labeled \textbf{GPIO2} is inserted in hole \textbf{D1} of the breadboard
    \item The pin labeled \textbf{GPIO20} is inserted in hole \textbf{H7} of the breadboard
    \item the pin labeled \textbf{GPIO21} is inserted in hole \textbf{D7} of the breadboard
\end{itemize}
You may need to apply more pressure than expected to seat the microcontroller properly in the breadboard. When its over, it should look like this:
\begin{figure}[H]
    \centering
    \includegraphics[width=.6\linewidth]{common/microcontroller_seated_in_breadboard.jpg}
    \caption{So far, so good!}
\end{figure}

\subsubsection{Connect the Display, Microphone, LED, and Resistor}
\begin{itemize}
    \item Place a red LED into the breadboard with the longer leg in \textbf{F20} and the shorter leg
    in \textbf{F21}. Then place one leg of a resistor in \textbf{J21} and the other in the top
    negative rail (the blue one).
    \item Place the 3 pins of the microphone board into the breadboard holes labeled \textbf{A12}, \textbf{A13}, and \textbf{A14}.
    \item Place the 4 pins of the OLED screen into the breadboard holes labeled \textbf{A26}, \textbf{A27}, \textbf{A28}, and \textbf{A29}.
\end{itemize}

It should now be looking like this:
\begin{figure}[H]
    \centering
    \includegraphics[width=.6\linewidth]{project_4/components_placed.jpg}
    \caption{All components are now ready to be wired up.}
\end{figure}

\subsubsection{Connect the necessary jumper wires}
\begin{itemize}
    \item Place one end of a black jumper wire into hole \textbf{J2} of the breadboard and the other end into
    the top negative rail (the blue one). This will provide a ground path for all of the components.
    \item Place a purple jumper wire between \textbf{J5} and \textbf{J20}. This provides the signal for the LED.
    \item Place a red jumper wire into hole \textbf{J3} of the breadboard and the other end in
    hole \textbf{C12} of the breadboard. This will provide \textbf{3.3} volts of power to the microphone.
    \item Place a black jumper wire into hole \textbf{C13} of the breadboard and the other end in the top negative rail.
    \item Place a yellow jumper wire into \textbf{C14} of the breadboard and the other end in hole \textbf{B3}.
    \item Place a red wire into \textbf{E12} and the other end into \textbf{D27} to provide power to the OLED.
    \item Place a brown wire into \textbf{E13} and the other end into \textbf{D26} to provide ground to the OLED.
    \item Place a white wire into \textbf{D28} and the other end into \textbf{B6} to provide a clock signal to the OLED.
    \item Place a yellow wire into \textbf{D29} and the other end into \textbf{B5} to provide data to the OLED.
\end{itemize}

You should be left with something that looks like this:
\begin{figure}[H]
    \centering
    \includegraphics[width=.6\linewidth]{project_4/all_wired_up.jpg}
    \caption{There power and ground wires are shared between the microphone and the OLED screen.}
\end{figure}

\subsection{Programming the microcontroller}
Once all of the wiring is correct, connect the USB cable to the microcontroller and load the IDE to
access it. Refer back to Chapter \ref{ide} for instructions.

Click on the file named "project\_4\_candle.py". This will load the code in the editor for this section.
Read through the comments and the code to get a sense for how it works. Once you are ready, you can
click the blue play button in the upper left of the window to start the script.

While the script is running, the LED will change brightness based on the sound level in the room. If
you blow on the microphone, the LED should get dimmer. If you blow long and hard enough, the LED will
go out. The OLED screen will show the current brightness level of the LED.

\section{Review}
This project demonstrates how you can connect external sensors to a microcontroller, read their status,
and output that status to a terminal or display. It also demonstrates how you can import external libraries
for talking to these peripheral devices.

\section{Possible Extensions}
If you want to do some experimentation, try these:

\begin{itemize}
    \item Connect an LED and have it light up when the temperature gets warmer than some threshold
    \item Connect a button and have the screen change between several different readouts when pressed
\end{itemize}
